Pressemitteilungen der Statistik Austria wie \glqq \acs{pkw}-Bestand 2025 um 54.208 höher als 2024\grqq{} oder \glqq 2025 so viele Pkw-Neuzulassungen wie seit 2019 nicht mehr\grqq{} verdeutlichen den anhaltenden Trend steigender Motorisierung in Österreich \parencites[S. 1]{statistikaustriaPkwBestand2025Um2026}[S. 1]{statistikaustria2025VielePkwNeuzulassungen2026}. 
Mit mehr als 7,5 Millionen zugelassenen \acp{kfz} wurde zwischen 1990 und 2025 ein Höchststand erreicht, wobei besonders große Fahrzeuge wie Geländewagen und \acp{suv} zunehmend öffentlichen Raum beanspruchen \parencite{hatzlDatenZeigenImmer2025}. 
Dieser Trend erhöht den Druck auf städtische Infrastrukturen und macht eine effizientere Nutzung vorhandener Verkehrsflächen erforderlich \parencite{jaschinskyWarumOsterreichsStadte2025}. 

Ein zentraler Aspekt dieser Entwicklung ist die zunehmende Schwierigkeit, geeignete Parkmöglichkeiten zu finden. 
Während Wien selbst einen vergleichsweise geringeren \acs{kfz}-Bestand aufweist, pendeln Erwerbstätige regelmäßig aus dem Umland in die Stadt \parencite{statistikaustriaKfzBestand2026}. 
Aufgrund teilweise unzureichender Anbindungen an den öffentlichen Verkehr erfolgt diese Mobilität häufig mit dem privaten \acs{kfz} \parencite{aknoAKNiederosterreichPendleranalyse20252025}. 

Die Suche nach einem Parkplatz endet dabei häufig in Parkhäusern, die nur teilweise mit Parkleitsystemen ausgestattet sind. 
Diese Systeme liefern meist lediglich statische oder lokal begrenzte Informationen, etwa zur Anzahl freier Stellplätze. 
Zwar bieten Funktionszonen in Wien eine grobe Orientierung, jedoch fehlt eine durchgängige, dynamische Unterstützung während der Anfahrt \parencite{stadtwienParkleitsystemFuerWien2026}. 

Digitale Navigationsdienste wie Google Maps unterstützen zwar die Routenplanung, integrieren jedoch keine aktuellen Informationen zur Parksituation. 
Nutzerinnen und Nutzer müssen daher auf unterschiedliche externe Informationsquellen zurückgreifen, was insbesondere während der Fahrt unpraktisch ist. 

Diese Situation wird zusätzlich durch das Ablenkungspotential mobiler Endgeräte verschärft. 
Das aktive Suchen und Verarbeiten von Informationen während der Fahrt beeinträchtigt nachweislich Fahrverhalten und Aufmerksamkeit \parencite[S. 361]{oviedo-trespalaciosUnderstandingImpactsMobile2016}. 
In Österreich ist rund ein Drittel der Verkehrsunfälle mit Todesfolge auf Unachtsamkeit und Ablenkung zurückzuführen \parencite{bmi397VerkehrstoteAuf2026}. 
Interaktionen wie Navigieren, Lesen oder Schreiben beeinflussen unter anderem Geschwindigkeit, Spurhaltung und Sicherheitsabstände negativ \parencite[S. 372-373]{oviedo-trespalaciosUnderstandingImpactsMobile2016}. 

Somit entsteht ein Spannungsfeld zwischen dem Bedarf an aktuellen Informationen und der sicheren Fahrzeugführung. 

Um diese Herausforderung zu adressieren, sollte die notwendige Informationsverarbeitung möglichst automatisiert und kontextabhängig erfolgen. 
Statt aktiv nach Alternativen zu suchen, könnten Nutzerinnen und Nutzer geeignete Optionen auf Basis aktueller Daten und individueller Präferenzen vorgeschlagen bekommen. 
Diese Vorschläge müssten während der Zielführung einfach bestätigt oder angepasst werden können, ohne die Aufmerksamkeit vom Verkehrsgeschehen abzulenken. 

Im Rahmen dieser Arbeit wird daher ein erster Lösungsansatz entwickelt, der eine dynamische Entscheidungsunterstützung während der Navigation ermöglicht. 
Durch den Einsatz von Sensoren zur Erfassung der Parkplatzsituation und die Integration in eine \ac{AWS}-basierte Infrastruktur wird die Grundlage für ein skalierbares \ac{IoT}-System geschaffen. 
Ziel ist es, relevante Informationen in Echtzeit bereitzustellen und gleichzeitig die Interaktion auf ein Minimum zu reduzieren.